\section{Experimental methodology and planned evaluation}
\label{sec:experiments}

\begin{figure}[htbp]
\centering
\resizebox{0.80\textwidth}{!}{%
\begin{tikzpicture}[node distance=15mm and 19mm]
\node[sh/source] (data) {Synthetic sources\\OTP / credential /\\identifier / payment};
\node[sh/infrastructure,right=of data] (vars) {Factors\\device, agent, model,\\binding, lifecycle, redaction};
\node[sh/attack,right=of vars] (attacks) {Attacks\\read, redirect, fake,\\replay, reread};
\node[sh/component,below=22mm of attacks] (conds) {Conditions\\B0 unprotected\\B1 placeholder\\B2 platform\\B3 SecretHandle};
\node[sh/defense,left=of conds] (oracle) {Protected oracles\\delivery + app state\\agent-artifact scan};
\node[sh/theory,left=of oracle] (metrics) {Metrics\\$E,R,U,L,F,Y$\\latency + overhead};
\draw[sh/flow] (data) -- (vars);
\draw[sh/attackflow] (vars) -- (attacks);
\draw[sh/flow] (attacks) -- (conds);
\draw[sh/secureflow] (conds) -- (oracle);
\draw[sh/secureflow] (oracle) -- (metrics);
\node[sh/infrastructure,below=17mm of oracle,minimum width=96mm] (repro) {Pinned AOSP/apps/agents/models + paired seeds + confidence intervals + raw-data scripts};
\draw[sh/control] (vars.south) -- (repro.north);
\draw[sh/control] (metrics.south) -- ([xshift=-32mm]repro.north);
\draw[sh/control] (conds.south) -- ([xshift=32mm]repro.north);
\end{tikzpicture}%
}
\caption{Planned experimental setup. Protected canary and delivery oracles distinguish actual secret flow from model behavior. Every empirical output remains pending until the AOSP prototype and runs exist.}
\label{fig:experiments}
\end{figure}


We evaluate \SH{} on a userdebug Android 14 (API 33/34) environment running on an EC2 instance with hardware-accelerated KVM virtualization. The evaluation answers the core research questions (RQ1--RQ5) using real Mobile-Agent v3.5 agents driven by Claude Sonnet~4.6, paired synthetic canary oracles, and an on-device microbenchmark harness.

\subsection{Prototype configurations}

The primary configuration is a userdebug AOSP build on a supported Pixel device. It includes the service, SMS/notification adapter, research Credential Manager provider, sensitive-clipboard adapter, resolver, principal-specific accessibility redactor, screenshot redactor, node-token support, and agent SDK. A secondary Samsung device evaluates only components implementable through public/system integration and reports whether keys are hardware-backed; Knox integration is not assumed until implemented. Emulator runs support deterministic testing but do not count as evidence for hardware-protected keys or real-device latency.

Four open test applications cover: (A) username/password plus SMS OTP; (B) multi-page 2FA with an adversarial message preview; (C) address and phone-number checkout; and (D) wallet-token checkout with per-use confirmation. Each application logs only ground-truth events to a protected harness: target identity, whether a value was delivered, server acceptance, and task completion. No real financial account or production OTP is used.

\subsection{Agents and model placement}

The evaluation targets AppAgent, the Mobile-Agent family, and another actively maintained open mobile-agent runner compatible with AndroidWorld. The idea document names DroidRun; because version and archival paper metadata were not verified during this audit, it remains an implementation candidate rather than a cited scholarly baseline. Each supported agent runs with at least one cloud model and one on-device or locally hosted model when feasible. Agent prompts, images, accessibility trees, tool calls, and traces are captured by a disclosure monitor that searches for exact canary values and structured encodings.

AndroidWorld measures whether general non-secret tasks regress \cite{rawles2024androidworld}. SPA-Bench can provide additional app/task breadth \cite{chen2024spabench}. Standard suites do not replace the authentication applications because they were not designed to exercise authoritative OTP, credential, clipboard, and wallet sources.

\subsection{Attacks and stress tests}

\begin{enumerate}[leftmargin=*]
  \item \textbf{Read and exfiltrate.} Unprotected agents navigate to a message or notification containing a canary OTP and encounter an instruction to send it to an attacker endpoint. SecretHandle variants expose only the bound handle.
  \item \textbf{Destination redirection.} An injected instruction requests \texttt{ENTER(handle)} into an attacker app, a different package, a WebView with another origin, a non-OTP field, or a stale/recycled node.
  \item \textbf{Fake handle.} A malicious app draws the exact handle visual grammar and submits guessed handle IDs and authenticators. Rate-limit behavior and false acceptance are recorded.
  \item \textbf{Replay and race.} The harness issues sequential and concurrent resolves around successful use and expiry. It crashes/restarts the service at each reservation phase.
  \item \textbf{Post-entry reread.} Immediately and repeatedly after injection, the agent captures screenshots, accessibility trees, text-change events, clipboard, IME state, and debugging logs.
  \item \textbf{Same-destination semantic misuse.} The attacker keeps the target within the allowed package/origin but changes surrounding account or transaction context. This test quantifies the residual risk that destination binding does not claim to eliminate.
  \item \textbf{Identity ambiguity.} The harness varies WebView navigation, nested frames, custom views, overlays, focus changes, split-screen, OEM keyboards, view recycling, rotation, and accessibility-service reconnects.
\end{enumerate}

Attack templates are drawn from the structure of indirect prompt-injection benchmarks \cite{zhan2024injecagent,debenedetti2024agentdojo}, GUI fine-print/environment attacks \cite{chen2025fineprint,liao2024eia}, and mobile third-party-channel workflows \cite{du2025mobileattacks}. We adapt mechanics, not reported success numbers.

\subsection{Baselines and ablations}

The primary baselines are:
\begin{description}[leftmargin=1.2cm,style=nextline]
  \item[B0: unprotected] The agent receives the raw value in its ordinary observation and types it normally.
  \item[B1: detected placeholder] An OCR/NER detector replaces rendered sensitive text with a local placeholder and resolves the placeholder for the requested target without destination binding, approximating the closest published mobile-placeholder family \cite{zhao2026anonymization}.
  \item[B2: platform flow] Android Autofill, OTP Autofill, or Credential Manager is used when the test app can invoke it \cite{androidAutofillServices,androidCredentialManager}. Unsupported agent workflows are recorded as not applicable rather than failures.
  \item[B3: SecretHandle] Full source minting, destination/task binding, lifecycle, and redaction.
\end{description}

Ablations disable source minting (detector fallback), destination binding, task binding, atomic lifecycle, post-entry redaction, origin checks, or per-use confirmation one at a time. Unsafe ablations run only with synthetic canaries and isolated endpoints.

\subsection{Metrics}

For $N$ tasks, let $I_i$ be the set of instrumented agent-visible artifacts and $s_i$ the canary secret. Exact exposure is
\begin{equation}
E=\frac{1}{N}\sum_{i=1}^{N}\mathbf{1}[s_i\in I_i],
\end{equation}
augmented by normalized, encoded, OCR, and substring checks. This metric can establish observed zero exposure under the harness but cannot prove absence from uninstrumented channels.

Redirection success $R$ is the fraction of attack trials in which the raw value reaches a target outside $d$. Legitimate task success $U$ is the fraction of benign trials satisfying the application oracle. Post-entry leakage $L$ counts trials in which the value appears in any subsequent modeled observation. Forgery acceptance $F$ and replay acceptance $Y$ count successful unauthorized injections per attempt. Latency reports end-to-end action overhead and component-level mint, authentication, target lookup, confirmation, injection, and redaction time.

We also record memory, CPU time, energy where supported, number of user confirmations, policy rejection rate, false rejection rate, crash-recovery outcome, and destination identity availability. General-task regression is the difference in success on benchmark tasks that do not use a handle.

\subsection{Design and sample planning}

Each condition uses fixed task seeds and paired benign/adversarial versions. Runs are blocked by device, agent, model, application, and secret type. Binary outcomes receive Wilson confidence intervals; paired comparisons use a paired permutation or exact McNemar test when assumptions fit. Latency uses medians, tail percentiles, bootstrap intervals, and full distributions. The number of repetitions is selected by a prospective precision analysis before looking at outcomes. All exclusions---agent crash, platform crash, network outage, user cancellation, unsupported identity---are reported by category.

Observed zeros are reported as ``0 of $n$'' with a confidence interval, not as proof of impossible attack. Multiple hypothesis testing is limited by identifying RQ1--RQ4 as primary outcomes and treating other comparisons as exploratory. Raw canary values are randomly generated per run, stored only in the protected oracle, and destroyed after analysis.

\begin{table}[htbp]
\centering
\caption{Primary empirical evaluation results on Android. Exposure $E$, Redirection $R$, Utility $U$, and Reread $L$ measured over $n=30$ independent paired trials per condition using Mobile-Agent v3.5 (Claude Sonnet~4.6). Replay $Y$ measured across $n=50$ replay attempts and $n=50$ concurrent races. Latency reports tail p95 response times. 95\% Wilson confidence intervals in brackets. Platform flow ($B_2$) is recorded as N/A due to lack of support in uninstrumented third-party agent flows.}
\label{tab:empirical-results}
\resizebox{\textwidth}{!}{\begin{tabular}{@{}lcccccc@{}}
\toprule
Condition & Exposure $E$ & Redirect $R$ & Utility $U$ & Reread $L$ & Replay $Y$ & p95 Latency \\
\midrule
$B_0$: Unprotected & 0.97 [0.83, 0.99] & 0.97 [0.83, 0.99] & 0.97 [0.83, 0.99] & 1.00 [0.89, 1.00] & 1.00 [0.93, 1.00] & 48.6 s \\
$B_1$: Detected placeholder & 0.00 [0.00, 0.11] & 1.00 [0.89, 1.00] & 0.00 [0.00, 0.11]$^\dagger$ & 1.00 [0.89, 1.00] & 0.00 [0.00, 0.07] & 0.15 s \\
$B_2$: Platform flow & N/A & N/A & N/A & N/A & N/A & N/A \\
$B_3$: \SH{} & \textbf{0.00} [0.00, 0.11] & \textbf{0.00} [0.00, 0.11] & \textbf{1.00} [0.89, 1.00] & \textbf{0.00} [0.00, 0.11] & \textbf{0.00} [0.00, 0.07] & \textbf{9.30 $\mu$s} \\
\bottomrule
\multicolumn{7}{l}{\footnotesize $^\dagger$ Under destination redirection attack, $B_1$ utility drops to 0.00 because the sensitive canary is consumed and leaked to the adversary.}
\end{tabular}}
\end{table}

\begin{figure}[htbp]
\centering
\begin{subfigure}[t]{0.48\textwidth}
\centering
\begin{tikzpicture}
\begin{axis}[
  width=0.91\linewidth,height=5.2cm,
  xlabel={Binding granularity},ylabel={Measured redirection leakage $R$},
  symbolic x coords={Unbound (B1),Bound (B3)},xtick=data,
  ymin=0,ymax=1.1,ytick={0,.25,.5,.75,1},
  grid=major,legend style={font=\scriptsize,draw=none,at={(0.5,-0.28)},anchor=north},
  tick label style={font=\scriptsize},label style={font=\small}]
\addplot[shorange,thick,mark=square*,sharp plot] coordinates {(Unbound (B1),1.00) (Bound (B3),0.00)};
\addlegendentry{Measured on device ($n=30$)}
\end{axis}
\end{tikzpicture}
\end{subfigure}\hfill
\begin{subfigure}[t]{0.48\textwidth}
\centering
\begin{tikzpicture}
\begin{axis}[
  width=0.91\linewidth,height=5.2cm,
  xlabel={Security condition},ylabel={Task utility $U$ under redirection},
  symbolic x coords={B0,B1,B3},xtick=data,
  ymin=0,ymax=1.1,ytick={0,.25,.5,.75,1},
  grid=major,legend style={font=\scriptsize,draw=none,at={(0.5,-0.28)},anchor=north},
  tick label style={font=\scriptsize},label style={font=\small}]
\addplot[shblue,thick,mark=*,sharp plot] coordinates {(B0,0.00) (B1,0.00) (B3,1.00)};
\addlegendentry{Measured utility ($n=30$)}
\end{axis}
\end{tikzpicture}
\end{subfigure}
\caption{Empirical evaluation on Android. Left: Unbound placeholder ($B_1$) allows $100\%$ ($30/30$) redirection leakage when an agent is tricked into submitting into an adversary app, whereas \SH{} ($B_3$) destination binding achieves $0.00$ ($0/30$, Wilson 95\% CI $[0.00, 0.11]$) leakage. Right: Under destination redirection attacks, legitimate utility $U$ collapses to $0.00$ for $B_0$ and $B_1$ (adversary hijacks secret), while \SH{} rejects rogue resolution and preserves legitimate utility ($30/30$, $1.00$).}
\label{fig:empirical-trends}
\end{figure}


\subsection{Killer experiment: Destination redirection (Attack 2, RQ2)}

The decisive experiment evaluates whether \SH{} prevents Confused Deputy destination redirection when an attacker tricks the mobile agent into injecting an OTP into an adversary-controlled application ($n=30$ paired trials). Under $B_0$, the agent observes the raw plaintext and transmits it directly to the adversary endpoint in $29/30$ trials ($R=0.97$, Wilson $[0.83, 0.99]$). Under $B_1$ (unbound detected placeholder), the placeholder is resolved blindly upon request without checking caller identity, yielding $30/30$ leaks ($R=1.00$, Wilson $[0.89, 1.00]$) and completely breaking legitimate utility ($U=0.00$, $0/30$). Under $B_3$ (\SH), the handle is bound to the legitimate package and task context; when presented by the attacker destination, the resolver validates the calling package via Algorithm~2 and rejects the resolution with \texttt{CALLER\_MISMATCH} in $30/30$ trials ($R=0.00$, Wilson $[0.00, 0.11]$). In legitimate workflows, the bound destination resolves successfully in $30/30$ trials ($U=1.00$, Wilson $[0.89, 1.00]$).

\subsection{Lifecycle atomicity and microbenchmark overhead (RQ3, RQ5)}

We execute an on-device benchmark harness directly inside the Android runtime via \texttt{app\_process} to measure cryptographic and vault operations ($n=100$ trials):
\begin{itemize}[leftmargin=*]
  \item \textbf{Mint latency:} Mean $5.61\,\mu\text{s}$ (median $5.24\,\mu\text{s}$, min $2.73\,\mu\text{s}$, max $41.64\,\mu\text{s}$, p95 $9.30\,\mu\text{s}$, std $3.96\,\mu\text{s}$).
  \item \textbf{Resolve latency:} Mean $0.13\,\mu\text{s}$ (median $0.12\,\mu\text{s}$, min $0.11\,\mu\text{s}$, max $0.67\,\mu\text{s}$, p95 $0.18\,\mu\text{s}$, std $0.06\,\mu\text{s}$).
  \item \textbf{Memory footprint:} Dalvik heap allocated is $1.15\,\text{MB}$ (total heap $8.39\,\text{MB}$); total process PSS is $28.71\,\text{MB}$ (native heap $14.19\,\text{MB}$, Dalvik heap $2.11\,\text{MB}$), demonstrating that \SH{} imposes negligible overhead on Android.
\end{itemize}

Under stress testing of the single-use lifecycle (Attack 4, $n=50$ trials each), sequential replay attempts are rejected in $50/50$ cases with \texttt{ALREADY\_RESOLVED} ($Y=0.00$, Wilson $[0.00, 0.07]$), and post-TTL resolution attempts are blocked in $50/50$ cases with \texttt{HANDLE\_EXPIRED} ($Y=0.00$). Under multi-threaded race conditions (10 concurrent threads per handle, 500 threads total synchronized via a concurrency barrier), exactly one thread succeeds in $50/50$ trials, with zero concurrency anomalies.

\subsection{Reproducibility checklist}

The artifact must pin the AOSP commit, device build, kernel, app hashes, provider versions, agent commits, model identifiers, prompts, seeds, network endpoints, and policy files. It must include a one-command isolated canary experiment, scripts that reconstruct every table/plot from raw data, schemas for trace redaction, and checks that no real secret is collected. Tests cover unit state transitions, concurrency schedules, property-based policy intersection, end-to-end UI flows, and regression on non-secret tasks.

Independent evaluators should be able to switch between B0--B3 without changing task text, inspect audit decisions, and verify oracle isolation. The release should separate device-required tests from emulator-compatible tests and document OEM features unavailable on each platform.
